%!TeX root=../treasuretop.tex
\ArtChapter[How My Sea Adventure Began]{chap22head}

   \lettrine[lines=4,image=true]{chap22initT}{here} was no return of the mutineers---not so much as another shot out of the woods. They had \enquote{got their rations for that day,} as the captain put it, and we had the place to ourselves and a quiet time to overhaul the wounded and get dinner. Squire and I cooked outside in spite of the danger, and even outside we could hardly tell what we were at, for horror of the loud groans that reached us from the doctor’s patients.

Out of the eight men who had fallen in the action, only three still breathed---that one of the pirates who had been shot at the loophole, Hunter, and Captain Smollett; and of these, the first two were as good as dead; the mutineer indeed died under the doctor’s knife, and Hunter, do what we could, never recovered consciousness in this world. He lingered all day, breathing loudly like the old buccaneer at home in his apoplectic fit, but the bones of his chest had been crushed by the blow and his skull fractured in falling, and some time in the following night, without sign or sound, he went to his Maker.

As for the captain, his wounds were grievous indeed, but not dangerous. No organ was fatally injured. Anderson’s ball---for it was Job that shot him first---had broken his shoulder-blade and touched the lung, not badly; the second had only torn and displaced some muscles in the calf. He was sure to recover, the doctor said, but in the meantime, and for weeks to come, he must not walk nor move his arm, nor so much as speak when he could help it.

My own accidental cut across the knuckles was a flea-bite. Doctor Livesey patched it up with plaster and pulled my ears for me into the bargain.

After dinner the squire and the doctor sat by the captain’s side awhile in consultation; and when they had talked to their hearts’ content, it being then a little past noon, the doctor took up his hat and pistols, girt on a cutlass, put the chart in his pocket, and with a musket over his shoulder crossed the palisade on the north side and set off briskly through the trees.

Gray and I were sitting together at the far end of the block house, to be out of earshot of our officers consulting; and Gray took his pipe out of his mouth and fairly forgot to put it back again, so thunder-struck he was at this occurrence.

\enquote{Why, in the name of Davy Jones,} said he, \enquote{is Dr Livesey mad?}

\enquote{Why no,} says I. \enquote{He’s about the last of this crew for that, I take it.}

\enquote{Well, shipmate,} said Gray, \enquote{mad he may not be; but if \textit{he’s} not, you mark my words, \textit{I} am.}

\enquote{I take it,} replied I, \enquote{the doctor has his idea; and if I am right, he’s going now to see Ben Gunn.}

I was right, as appeared later; but in the meantime, the house being stifling hot and the little patch of sand inside the palisade ablaze with midday sun, I began to get another thought into my head, which was not by any means so right. What I began to do was to envy the doctor walking in the cool shadow of the woods with the birds about him and the pleasant smell of the pines, while I sat grilling, with my clothes stuck to the hot resin, and so much blood about me and so many poor dead bodies lying all around that I took a disgust of the place that was almost as strong as fear.

All the time I was washing out the block house, and then washing up the things from dinner, this disgust and envy kept growing stronger and stronger, till at last, being near a bread-bag, and no one then observing me, I took the first step towards my escapade and filled both pockets of my coat with biscuit.

I was a fool, if you like, and certainly I was going to do a foolish, over-bold act; but I was determined to do it with all the precautions in my power. These biscuits, should anything befall me, would keep me, at least, from starving till far on in the next day.

The next thing I laid hold of was a brace of pistols, and as I already had a powder-horn and bullets, I felt myself well supplied with arms.

As for the scheme I had in my head, it was not a bad one in itself. I was to go down the sandy spit that divides the anchorage on the east from the open sea, find the white rock I had observed last evening, and ascertain whether it was there or not that Ben Gunn had hidden his boat, a thing quite worth doing, as I still believe. But as I was certain I should not be allowed to leave the enclosure, my only plan was to take French leave and slip out when nobody was watching, and that was so bad a way of doing it as made the thing itself wrong. But I was only a boy, and I had made my mind up.

Well, as things at last fell out, I found an admirable opportunity. The squire and Gray were busy helping the captain with his bandages, the coast was clear, I made a bolt for it over the stockade and into the thickest of the trees, and before my absence was observed I was out of cry of my companions.

This was my second folly, far worse than the first, as I left but two sound men to guard the house; but like the first, it was a help towards saving all of us.

I took my way straight for the east coast of the island, for I was determined to go down the sea side of the spit to avoid all chance of observation from the anchorage. It was already late in the afternoon, although still warm and sunny. As I continued to thread the tall woods, I could hear from far before me not only the continuous thunder of the surf, but a certain tossing of foliage and grinding of boughs which showed me the sea breeze had set in higher than usual. Soon cool draughts of air began to reach me, and a few steps farther I came forth into the open borders of the grove, and saw the sea lying blue and sunny to the horizon and the surf tumbling and tossing its foam along the beach.

I have never seen the sea quiet round Treasure Island. The sun might blaze overhead, the air be without a breath, the surface smooth and blue, but still these great rollers would be running along all the external coast, thundering and thundering by day and night; and I scarce believe there is one spot in the island where a man would be out of earshot of their noise.

I walked along beside the surf with great enjoyment, till, thinking I was now got far enough to the south, I took the cover of some thick bushes and crept warily up to the ridge of the spit.

Behind me was the sea, in front the anchorage. The sea breeze, as though it had the sooner blown itself out by its unusual violence, was already at an end; it had been succeeded by light, variable airs from the south and south-east, carrying great banks of fog; and the anchorage, under lee of Skeleton Island, lay still and leaden as when first we entered it. The \textit{Hispaniola}, in that unbroken mirror, was exactly portrayed from the truck to the waterline, the Jolly Roger hanging from her peak.

Alongside lay one of the gigs, Silver in the stern-sheets---him I could always recognize---while a couple of men were leaning over the stern bulwarks, one of them with a red cap---the very rogue that I had seen some hours before stride-legs upon the palisade. Apparently they were talking and laughing, though at that distance---upwards of a mile---I could, of course, hear no word of what was said. All at once there began the most horrid, unearthly screaming, which at first startled me badly, though I had soon remembered the voice of Captain Flint and even thought I could make out the bird by her bright plumage as she sat perched upon her master’s wrist.

Soon after, the jolly-boat shoved off and pulled for shore, and the man with the red cap and his comrade went below by the cabin companion.

Just about the same time, the sun had gone down behind the Spy-glass, and as the fog was collecting rapidly, it began to grow dark in earnest. I saw I must lose no time if I were to find the boat that evening.

The white rock, visible enough above the brush, was still some eighth of a mile further down the spit, and it took me a goodish while to get up with it, crawling, often on all fours, among the scrub. Night had almost come when I laid my hand on its rough sides. Right below it there was an exceedingly small hollow of green turf, hidden by banks and a thick underwood about knee-deep, that grew there very plentifully; and in the centre of the dell, sure enough, a little tent of goat-skins, like what the gipsies carry about with them in England.

I dropped into the hollow, lifted the side of the tent, and there was Ben Gunn’s boat---home-made if ever anything was home-made; a rude, lop-sided framework of tough wood, and stretched upon that a covering of goat-skin, with the hair inside. The thing was extremely small, even for me, and I can hardly imagine that it could have floated with a full-sized man. There was one thwart set as low as possible, a kind of stretcher in the bows, and a double paddle for propulsion.

  \begin{sidewaysfigure}
\includegraphics[width=\columnwidth]{chap22bengunnsboat}%
\caption{There was Ben Gunn’s boat}
\end{sidewaysfigure} 
 

I had not then seen a coracle, such as the ancient Britons made, but I have seen one since, and I can give you no fairer idea of Ben Gunn’s boat than by saying it was like the first and the worst coracle ever made by man. But the great advantage of the coracle it certainly possessed, for it was exceedingly light and portable.

Well, now that I had found the boat, you would have thought I had had enough of truantry for once, but in the meantime I had taken another notion and become so obstinately fond of it that I would have carried it out, I believe, in the teeth of Captain Smollett himself. This was to slip out under cover of the night, cut the \textit{Hispaniola} adrift, and let her go ashore where she fancied. I had quite made up my mind that the mutineers, after their repulse of the morning, had nothing nearer their hearts than to up anchor and away to sea; this, I thought, it would be a fine thing to prevent, and now that I had seen how they left their watchmen unprovided with a boat, I thought it might be done with little risk.

Down I sat to wait for darkness, and made a hearty meal of biscuit. It was a night out of ten thousand for my purpose. The fog had now buried all heaven. As the last rays of daylight dwindled and disappeared, absolute blackness settled down on Treasure Island. And when, at last, I shouldered the coracle and groped my way stumblingly out of the hollow where I had supped, there were but two points visible on the whole anchorage.

One was the great fire on shore, by which the defeated pirates lay carousing in the swamp. The other, a mere blur of light upon the darkness, indicated the position of the anchored ship. She had swung round to the ebb---her bow was now towards me---the only lights on board were in the cabin, and what I saw was merely a reflection on the fog of the strong rays that flowed from the stern window.

The ebb had already run some time, and I had to wade through a long belt of swampy sand, where I sank several times above the ankle, before I came to the edge of the retreating water, and wading a little way in, with some strength and dexterity, set my coracle, keel downwards, on the surface.